\documentclass[11pt]{article}
\usepackage{amsmath,amssymb,amsthm}
\usepackage{algorithm}
\usepackage{algpseudocode}
\usepackage{fullpage}
\usepackage{hyperref}
\newtheorem{theorem}{Theorem}
\newtheorem{lemma}{Lemma}
\newtheorem{assumption}{Assumption}
\newcommand{\E}{\mathbb{E}}
\newcommand{\R}{\mathbb{R}}

\begin{document}

\section*{Algorithm Name}
\textbf{SAGA (Stochastic Average Gradient Augmented)} – a variance‑reduced stochastic gradient method that stores a table of past component gradients and uses them to construct an unbiased gradient estimator.

\section*{Mathematical Formulation}
Consider the finite‑sum convex optimization problem
\[
\min_{x\in\R^{d}} \; f(x) \;:=\; \frac{1}{n}\sum_{i=1}^{n} f_i(x),
\]
where each $f_i:\R^{d}\!\to\!\R$ is convex and $L$‑smooth.  
Let $\phi_i^{k}\in\R^{d}$ denote the point at which the $i$‑th component gradient was last stored (initially $\phi_i^{0}=x^{0}$).  
Define the \emph{gradient table}
\[
g_i^{k}\;:=\;\nabla f_i(\phi_i^{k}),\qquad 
\bar g^{k}\;:=\;\frac{1}{n}\sum_{i=1}^{n} g_i^{k}.
\]

At iteration $k$, draw an index $i_k\in\{1,\dots,n\}$ uniformly at random and compute the \emph{SAGA estimator}
\[
\tilde{\nabla}_k \;:=\; \nabla f_{i_k}(x^{k}) - g_{i_k}^{k} + \bar g^{k}.
\]
The update rule is
\[
\boxed{\,x^{k+1}=x^{k}-\alpha \tilde{\nabla}_k\,},
\]
where $\alpha>0$ is the stepsize.  
After the update, replace the stored gradient for the sampled component:
\[
\phi_{i_k}^{k+1}=x^{k},\qquad g_{i_k}^{k+1}= \nabla f_{i_k}(x^{k}),
\]
and keep all other table entries unchanged:
\[
\phi_{j}^{k+1}= \phi_{j}^{k},\; g_{j}^{k+1}=g_{j}^{k},\;\forall j\neq i_k .
\]

\section*{Pseudocode}
\begin{algorithm}[H]
\caption{SAGA}
\begin{algorithmic}[1]
\State \textbf{Input:} stepsize $\alpha>0$, number of epochs $T$.
\State Initialize $x^{0}\in\R^{d}$.
\For{$i=1,\dots,n$}
    \State $\phi_i^{0}\gets x^{0}$,\quad $g_i^{0}\gets \nabla f_i(x^{0})$
\EndFor
\State $\bar g^{0}\gets \frac{1}{n}\sum_{i=1}^{n} g_i^{0}$
\For{$k=0,1,\dots,T-1$}
    \State Sample $i_k\sim\text{Uniform}\{1,\dots,n\}$
    \State $\tilde{\nabla}_k \gets \nabla f_{i_k}(x^{k})-g_{i_k}^{k}+\bar g^{k}$
    \State $x^{k+1}\gets x^{k}-\alpha \tilde{\nabla}_k$
    \State $g_{i_k}^{k+1}\gets \nabla f_{i_k}(x^{k})$
    \State $\phi_{i_k}^{k+1}\gets x^{k}$
    \State $\bar g^{k+1}\gets \bar g^{k}+ \frac{1}{n}\bigl(g_{i_k}^{k+1}-g_{i_k}^{k}\bigr)$
\EndFor
\State \textbf{Output:} $x^{T}$
\end{algorithmic}
\end{algorithm}

\section*{Dry Run Example}
Minimize the scalar quadratic
\[
f(w)=\frac{1}{2}(w-3)^{2},
\]
which can be written as a finite sum with $n=1$: $f_1(w)=\frac{1}{2}(w-3)^{2}$.
Thus $g_1^{k}= \nabla f_1(\phi_1^{k}) = \phi_1^{k}-3$ and $\bar g^{k}=g_1^{k}$.

\begin{center}
\begin{tabular}{c|c|c|c}
iteration $k$ & $x^{k}$ & $\tilde{\nabla}_k$ & $f(x^{k})$\\\hline
0 & $0$ & $\bigl[(0-3)- (0-3)+(0-3)\bigr]= -3$ & $4.5$\\
1 & $0-0.1(-3)=0.3$ & $[(0.3-3)-(0-3)+(-3)] = -2.7$ & $3.645$\\
2 & $0.3-0.1(-2.7)=0.57$ & $[(0.57-3)- (0.3-3)+(-2.7)] = -2.13$ & $2.9529$\\
3 & $0.57-0.1(-2.13)=0.783$ & $[(0.783-3)-(0.57-3)+(-2.13)] = -1.767$ & $2.3717$
\end{tabular}
\end{center}
The objective strictly decreases, illustrating the variance‑reduced nature of the estimator.

\newpage
\section*{Assumptions}
\begin{assumption}[Convexity]\label{ass:convex}
Each component $f_i$ is convex:
\[
f_i(y)\ge f_i(x)+\langle \nabla f_i(x),y-x\rangle,\qquad \forall x,y\in\R^{d}.
\]
\end{assumption}
\begin{assumption}[Smoothness]\label{ass:smooth}
Each $f_i$ is $L$‑smooth:
\[
\|\nabla f_i(y)-\nabla f_i(x)\|\le L\|y-x\|,\qquad \forall x,y\in\R^{d}.
\]
\end{assumption}
\begin{assumption}[Finite Sum]\label{ass:finitesum}
The objective is the average of $n$ components:
\[
f(x)=\frac{1}{n}\sum_{i=1}^{n}f_i(x).
\]
\end{assumption}
\begin{assumption}[Stepsize]\label{ass:stepsize}
The stepsize satisfies $0<\alpha\le \frac{1}{3L}$.
\end{assumption}

\section*{Main Convergence Theorem}
\begin{theorem}\label{thm:linear}
Assume \ref{ass:convex}–\ref{ass:stepsize} and that $f$ is $\mu$‑strongly convex ($\mu>0$).  
Define the Lyapunov function
\[
\Phi^{k}\;:=\;\|x^{k}-x^{*}\|^{2}+ \frac{\alpha}{n}\sum_{i=1}^{n}\|\phi_i^{k}-x^{*}\|^{2},
\]
where $x^{*}$ is the unique minimizer of $f$.  
Then for all $k\ge0$,
\[
\boxed{\E[\Phi^{k+1}] \;\le\; (1-\alpha\mu)\,\E[\Phi^{k}]}.
\]
Consequently,
\[
\E\bigl[f(x^{k})-f(x^{*})\bigr]\;\le\; \frac{L}{2}\,(1-\alpha\mu)^{k}\,\Phi^{0}.
\]
Thus SAGA attains a linear (geometric) convergence rate.
\end{theorem}
\begin{theorem}\label{thm:convex}
If $f$ is merely convex (no strong convexity) and the stepsize obeys $\alpha\le \frac{1}{3L}$, then for the averaged iterate $\bar x^{K}:=\frac{1}{K}\sum_{k=0}^{K-1}x^{k}$,
\[
\E\bigl[f(\bar x^{K})-f(x^{*})\bigr]\;\le\; \frac{2L\Phi^{0}}{K}.
\]
Hence SAGA achieves the optimal $O(1/K)$ rate for convex finite‑sum problems.
\end{theorem}

\section*{Theoretical Properties}
\begin{itemize}
\item \textbf{Unbiasedness.} The estimator $\tilde{\nabla}_k$ satisfies $\E[\tilde{\nabla}_k\mid x^{k}]=\nabla f(x^{k})$.
\item \textbf{Variance Reduction.} The variance of $\tilde{\nabla}_k$ decays as the stored points $\phi_i^{k}$ approach $x^{*}$; formally,
\[
\E\!\bigl[\|\tilde{\nabla}_k-\nabla f(x^{k})\|^{2}\bigr]\le
\frac{L^{2}}{n}\sum_{i=1}^{n}\E\!\bigl[\|\phi_i^{k}-x^{*}\|^{2}\bigr].
\]
\item \textbf{Stepsize Sensitivity.} The linear rate requires $\alpha\le 1/(3L)$; larger stepsizes break the contractive property of the Lyapunov function.
\item \textbf{Comparison.} Compared with vanilla SGD (which attains $O(1/\sqrt{K})$ under convexity), SAGA enjoys $O(1/K)$ without diminishing stepsizes, and a geometric rate under strong convexity, matching the optimal rates of variance‑reduced methods such as SVRG and SAG.
\end{itemize}

\section*{Discussion}
The theorem shows that by maintaining a gradient table, SAGA eliminates the stochastic gradient noise asymptotically. The linear convergence constant $(1-\alpha\mu)$ is identical to that of full‑gradient descent with stepsize $\alpha$, yet each iteration costs only a single component gradient evaluation. The $O(1/K)$ guarantee in the merely convex case demonstrates that diminishing stepsizes are unnecessary, simplifying practical tuning.

\newpage
\section*{Preliminaries (Definitions and Lemmas)}
\begin{lemma}[Smoothness Inequality]\label{lem:smooth}
If $f_i$ is $L$‑smooth then for all $x,y$,
\[
f_i(y)\le f_i(x)+\langle\nabla f_i(x),y-x\rangle+\frac{L}{2}\|y-x\|^{2}.
\]
\end{lemma}
\begin{lemma}[Strong Convexity Inequality]\label{lem:strong}
If $f$ is $\mu$‑strongly convex then for all $x$,
\[
\| \nabla f(x) \|^{2}\ge 2\mu\bigl[f(x)-f(x^{*})\bigr].
\]
\end{lemma}
Both lemmas follow from standard convex analysis (see e.g., Nesterov, 2004).

\section*{Main Proof (Step‑by‑Step Derivation)}
We prove Theorem~\ref{thm:linear}.  
All expectations are taken w.r.t. the random index $i_k$ conditioned on the past.

\noindent\textbf{Step 1.} Expand the distance term.  
\begin{align}
\|x^{k+1}-x^{*}\|^{2}
&= \bigl\|x^{k}-\alpha\tilde{\nabla}_k-x^{*}\bigr\|^{2}\nonumber\\
&= \|x^{k}-x^{*}\|^{2}
-2\alpha\langle\tilde{\nabla}_k,\,x^{k}-x^{*}\rangle
+\alpha^{2}\|\tilde{\nabla}_k\|^{2}. \tag{Step 1}
\end{align}

\noindent\textbf{Step 2.} Take conditional expectation and use unbiasedness.  
Because $\E[\tilde{\nabla}_k\mid x^{k}]=\nabla f(x^{k})$, we have
\[
\E\bigl[\langle\tilde{\nabla}_k,x^{k}-x^{*}\rangle\mid x^{k}\bigr]
= \langle\nabla f(x^{k}),x^{k}-x^{*}\rangle. \tag{Step 2}
\]

\noindent\textbf{Step 3.} Bound the second‑moment term.  
Write $\tilde{\nabla}_k$ as
\[
\tilde{\nabla}_k = \nabla f_{i_k}(x^{k})-g_{i_k}^{k}+\bar g^{k}.
\]
Using $g_{i_k}^{k}= \nabla f_{i_k}(\phi_{i_k}^{k})$, we obtain
\[
\tilde{\nabla}_k-\nabla f(x^{k}) = \bigl[\nabla f_{i_k}(x^{k})-\nabla f_{i_k}(\phi_{i_k}^{k})\bigr]
-\bigl[\bar g^{k}-\nabla f(x^{k})\bigr].
\]
Applying the inequality $\|a+b\|^{2}\le 2\|a\|^{2}+2\|b\|^{2}$ and smoothness,
\begin{align}
\E\bigl[\|\tilde{\nabla}_k\|^{2}\mid x^{k}\bigr]
&\le 2\E\bigl[\|\nabla f_{i_k}(x^{k})\|^{2}\bigr]+2\|\bar g^{k}\|^{2}\nonumber\\
&\le 4L^{2}\Bigl(\|x^{k}-x^{*}\|^{2}
+\frac{1}{n}\sum_{i=1}^{n}\|\phi_i^{k}-x^{*}\|^{2}\Bigr). \tag{Step 3}
\end{align}
(The detailed algebra: by $L$‑smoothness,
$\|\nabla f_i(x^{k})-\nabla f_i(\phi_i^{k})\|\le L\|x^{k}-\phi_i^{k}\|$;
then $\|x^{k}-\phi_i^{k}\|^{2}\le 2\|x^{k}-x^{*}\|^{2}+2\|\phi_i^{k}-x^{*}\|^{2}$.)

\noindent\textbf{Step 4.} Use convexity to bound the inner product.  
From convexity of $f$,
\[
\langle\nabla f(x^{k}),x^{k}-x^{*}\rangle\ge f(x^{k})-f(x^{*}). \tag{Step 4}
\]

\noindent\textbf{Step 5.} Combine Steps 1–4.  
Taking full expectation and substituting the bounds,
\begin{align}
\E\bigl[\|x^{k+1}-x^{*}\|^{2}\bigr]
&\le \E\bigl[\|x^{k}-x^{*}\|^{2}\bigr]
-2\alpha\,\E\bigl[f(x^{k})-f(x^{*})\bigr]
+2\alpha^{2}L^{2}\E\!\Bigl[\|x^{k}-x^{*}\|^{2}
+\frac{1}{n}\sum_{i=1}^{n}\|\phi_i^{k}-x^{*}\|^{2}\Bigr]. \tag{Step 5}
\end{align}

\noindent\textbf{Step 6.} Update the table term.  
Only the sampled index $i_k$ changes:
\[
\|\phi_{i_k}^{k+1}-x^{*}\|^{2}= \|x^{k}-x^{*}\|^{2},
\qquad
\|\phi_{j}^{k+1}-x^{*}\|^{2}= \|\phi_{j}^{k}-x^{*}\|^{2},\;j\neq i_k.
\]
Hence, taking expectation over $i_k$,
\begin{align}
\E\!\Bigl[\frac{1}{n}\sum_{i=1}^{n}\|\phi_i^{k+1}-x^{*}\|^{2}\Bigr]
&= \frac{1}{n}\E\!\bigl[\|x^{k}-x^{*}\|^{2}\bigr]
+ \Bigl(1-\frac{1}{n}\Bigr)\frac{1}{n}\sum_{i=1}^{n}\E\!\bigl[\|\phi_i^{k}-x^{*}\|^{2}\bigr]. \tag{Step 6}
\end{align}

\noindent\textbf{Step 7.} Form the Lyapunov function.  
Define
\[
\Phi^{k}= \|x^{k}-x^{*}\|^{2}+ \alpha\frac{1}{n}\sum_{i=1}^{n}\|\phi_i^{k}-x^{*}\|^{2}.
\]
Add $\alpha$ times the relation in Step 6 to the inequality in Step 5:
\begin{align}
\E[\Phi^{k+1}]
&\le \bigl(1-2\alpha\mu +2\alpha^{2}L^{2}\bigr)\E[\|x^{k}-x^{*}\|^{2}]
+ \bigl(1-\tfrac{1}{n}+2\alpha^{2}L^{2}\bigr)\alpha\frac{1}{n}\sum_{i=1}^{n}\E[\|\phi_i^{k}-x^{*}\|^{2}]\nonumber\\
&\qquad -2\alpha\,\E[f(x^{k})-f(x^{*})]. \tag{Step 7}
\end{align}
Here we used Lemma~\ref{lem:strong} to replace $f(x^{k})-f(x^{*})\ge \frac{\mu}{2}\|x^{k}-x^{*}\|^{2}$.

\noindent\textbf{Step 8.} Choose $\alpha\le\frac{1}{3L}$ and use $\mu\le L$.  
Under this stepsize,
\[
1-2\alpha\mu +2\alpha^{2}L^{2}\le 1-\alpha\mu,\qquad
1-\frac{1}{n}+2\alpha^{2}L^{2}\le 1-\frac{1}{n}+ \frac{2}{9}\le 1.
\]
Thus (Step 7) simplifies to
\[
\E[\Phi^{k+1}] \le (1-\alpha\mu)\,\E[\Phi^{k}]. \tag{Step 8}
\]

\noindent\textbf{Step 9.} Recurrence solution.  
Iterating the contraction yields
\[
\E[\Phi^{k}] \le (1-\alpha\mu)^{k}\Phi^{0}. \tag{Step 9}
\]

\noindent\textbf{Step 10.} Relate $\Phi^{k}$ to function suboptimality.  
From smoothness,
\[
f(x^{k})-f(x^{*})\le \frac{L}{2}\|x^{k}-x^{*}\|^{2}\le \frac{L}{2}\Phi^{k}.
\]
Taking expectations and using (Step 9) gives the claimed bound in Theorem~\ref{thm:linear}. \hfill $\square$

\section*{Rate Derivation (With Constants)}
From Theorem~\ref{thm:linear}:
\[
\E\bigl[f(x^{k})-f(x^{*})\bigr]
\le \frac{L}{2}(1-\alpha\mu)^{k}\Phi^{0},
\qquad 
\Phi^{0}= \|x^{0}-x^{*}\|^{2}+ \alpha\frac{1}{n}\sum_{i=1}^{n}\|\phi_i^{0}-x^{*}\|^{2}.
\]
Choosing the maximal admissible stepsize $\alpha=1/(3L)$ maximizes the contraction factor:
\[
1-\alpha\mu = 1-\frac{\mu}{3L}.
\]
Hence the linear rate constant is $1-\mu/(3L)$.

For the merely convex case (Theorem~\ref{thm:convex}), a telescoping sum applied to the inequality obtained from Step 5 (without strong convexity) yields
\[
\frac{1}{K}\sum_{k=0}^{K-1}\E\bigl[f(x^{k})-f(x^{*})\bigr]
\le \frac{2L\Phi^{0}}{K},
\]
which after Jensen’s inequality gives the $O(1/K)$ bound for the averaged iterate.

\section*{Special Cases}
\begin{itemize}
\item \textbf{Quadratic objectives.} When each $f_i(x)=\frac{1}{2}x^{\top}A_i x-b_i^{\top}x$, the SAGA estimator becomes an exact linear combination of the stored gradients, and the contraction factor can be expressed in terms of the smallest eigenvalue of $\frac{1}{n}\sum_i A_i$.
\item \textbf{Infinite data limit.} If $n\to\infty$ and the table is initialized with a single sample, the method reduces to standard SGD because the average table gradient $\bar g^{k}$ vanishes in expectation. The convergence reverts to $O(1/\sqrt{K})$.
\item \textbf{Non‑uniform sampling.} The proof extends to sampling probabilities $p_i$ with minor modifications: the stepsize condition becomes $\alpha\le \frac{1}{3L_{\max}}$ where $L_{\max}=\max_i L_i/p_i$.
\end{itemize}

\end{document}